% Generated draft: requires mathematical and source review.
\documentclass[11pt]{article}
\usepackage[margin=25mm]{geometry}
\usepackage{fontspec}
\setmainfont{Latin Modern Roman}
\usepackage{amsmath,amssymb,mathtools,unicode-math}
\setmathfont{Latin Modern Math}
\usepackage{enumitem,needspace,xurl,hyperref}
\setlength{\parindent}{0pt}
\setlength{\parskip}{5pt}
\setlength{\emergencystretch}{3em}
\newcommand{\N}{\mathbb{N}}
\newcommand{\Z}{\mathbb{Z}}
\newcommand{\Q}{\mathbb{Q}}
\newcommand{\R}{\mathbb{R}}
\newcommand{\C}{\mathbb{C}}
\newcommand{\F}{\mathbb{F}}
\newcommand{\A}{\mathbb{A}}
\newcommand{\PP}{\mathbb{P}}
\newcommand{\E}{\mathbb{E}}
\newcommand{\eps}{\varepsilon}
\newcommand{\dd}{\,\mathrm{d}}
\begin{document}
\section*{UnsolvedMath ID: 30004581}
\textbf{Source problem code:} catalogue entry
\subsection{Definitions and mathematical statement}
The supplied record poses a historical inquiry rather than a mathematical conjecture. It supplies no operational definitions of mathematical practitioners, professional status, comparison disciplines, public attitudes, or effects on a profession. The geographical scope and chronological boundaries of the medieval period are also unspecified. Consequently, no precise mathematical proposition or system of quantifiers can be extracted without adding assumptions.

The source question is retained verbatim:
\begin{quote}
In the medieval period, what was mathematical practitioners' professional status relative to other disciplines, and how did public attitudes affect their profession?
\end{quote}

A precise historical formulation would require specified periods, regions, practitioner groups, comparison disciplines, documentary evidence, and criteria for assessing professional status and the influence of public attitudes. These remain unresolved rather than being supplied here. The registry labels the question open; this is a supplied classification, not a verified mathematical fact.
\subsection{Short English statement}
The question asks how medieval mathematical practitioners ranked professionally compared with practitioners of other disciplines, and how public attitudes influenced their profession. Its historical scope and evaluative terms require clarification before a precise research target can be stated; the supplied record does not determine a mathematical conjecture.
\subsection{Sources}
\begin{itemize}
\item Local UnsolvedMath registry, permanent ID \(30004581\), title: Professional Status of Medieval Mathematical Practitioners; source problem code \texttt{catalogue entry}. The supplied record classifies the statement as exact and reports that it was internally readable, complete, and retained byte-for-byte. These are supplied textual qualifications, not mathematical verification.
\item Supplied source citation: \emph{History of Mathematics: A Global Cultural Approach} (2021), Oberwolfach Reports, pp. 2021--2092, DOI 10.4171/owr/2020/41. \url{https://doi.org/10.4171/owr/2020/41}. The background describes Oberwolfach Reports as documenting mathematical workshops and their open-problem discussions. The registry gives proposed year 2020 and identifies no individual proposer.
\item Supplied literature source, labelled ``original Oberwolfach report'': \url{https://doi.org/10.4171/owr/2020/41}.
\item Literature assessment dated exactly \texttt{2026-08-22}: the supplied curation labels the inquiry open and describes it as a broad historical inquiry with no uniquely verifiable final answer in the cited literature. This is supplied dated curation, not an independently verified assessment. The source background requests independent verification of time-sensitive status and cited claims.
\end{itemize}
\end{document}
